\section{Reproducibility statement}
\label{sec:repro}

\paragraph{Data.} The benchmark is generated, not distributed as a fixed file: seed \nSeed{} and
the generator code at a stated commit reproduce a draw byte for byte, independently of how the work
is batched. Each draw is identified by a SHA-256 fingerprint over its output rows (\texttt{6abde44e}
and \texttt{d8083dbc} in this paper), and its realism report records the parameter hash, the check
set hash and the fingerprint. The dataset release is licensed CC BY 4.0
(\texttt{dataset/release/}); the code is Apache-2.0.

\paragraph{Evidence.} Every figure in this paper comes from a raw output file (Appendix~\ref{app:evidence}) under
\texttt{docs/benchmarks/} whose header records the command, its exit status, the commit and whether
the working tree was clean; \texttt{numbers.tex} maps each figure in the text to its file. Every
read of the test period is logged in \texttt{docs/benchmarks/test\_set\_access.jsonl}. Experiment
tracking in MLflow was not used for the evidence in this paper (ADR 0030); the evidence files carry the
provenance instead.

\paragraph{Hardware.} All measurements were made on one development laptop (\nHwCPU{},
\nHwMem{}), recorded in \texttt{docs/benchmarks/hardware.md}. Latency figures from it are reported
as shapes, not as gate numbers (ADR 0010).

\paragraph{One command.} A single target that regenerates the dataset, trains, evaluates and
produces the paper's tables and figures from a clean machine, with compute time and hardware
recorded, is \notyet{M11}. Until it exists, each figure is reproduced by the command recorded in its
evidence file.

\paragraph{Publication.} Scripts and checklists for publishing the dataset and model, for the author
to run, are part of this milestone and do not exist yet. No DOI exists, and none is cited.
